% Section content template for individual Educative lessons
% This file contains the actual content for a specific section
% Generated from Educative API section content

% Section: Summary: Designing a Vending Machine
% Section ID: 4847744110690304
% Section Slug: summary-designing-a-vending-machine
% Generated: 2025-12-11T20:26:26.581035

Now that you've completed the vending machine case study, let's take a moment to reflect on and consolidate what we've learned. We 'll revisit the key system requirements, identify the core classes along with their responsibilities and relationships, and highlight the major design principles applied. We
'll also examine how objects interact within the system and walk through the overall workflow to understand how the components come together to achieve the desired functionality.

\subsection{Key requirements}\label{dPBZqYkhZOn4JeAD6dU6v}
\\
This section outlines the primary functional requirements that shaped the design of the vending machine system.

\begin{enumerate}
\item
\protect\phantomsection\label{f7Y9af86sj37-Jhyhp2T1}
The machine must store and manage various products, assigning each item to a unique slot.
\item
\protect\phantomsection\label{_PpeaumDEb_Uh7iJ-X9pn}
The system operates in three distinct states: \texttt{NoMoneyInsertedState}, \texttt{MoneyInsertedState}, and \texttt{DispenseState}.
\item
\protect\phantomsection\label{XhKKrUwMEIk1jkB4Yte_L}
The system has two main actors: a Customer, who buys products, and an Operator, who manages the machine.
\item
\protect\phantomsection\label{0TVOimskkHjQE3sfbDxnX}
An Operator must be able to add or remove products from the machine's inventory.
\item
\protect\phantomsection\label{jqB3qBtjs6_cJ54AH_Fnx}
Users must be able to select a product by specifying its rack number.
\item
\protect\phantomsection\label{pUlRJdoTvGfYUmd9C7LPQ}
Users can insert cash as payment.
\item
\protect\phantomsection\label{n82ixOi4Lge8JgTDY91oX}
The system must accurately calculate the total amount of money inserted.
\item
\protect\phantomsection\label{P1nUaRaxGrYBj-lfda-y7}
The system must verify if the inserted amount is sufficient for the selected product.
\item
\protect\phantomsection\label{hVO3F6saz8ObU_BJpmXha}
If the user pays more than the product price, the system must dispense the product and return the correct change.
\item
\protect\phantomsection\label{pqYYruerLSQEY1uAA0p5f}
If the user pays less than the product price, the system must display an error and return the inserted money.
\end{enumerate}

\subsection{System actors}\label{C1bBcCIycbj5FasxDKIgj}
\\
The following actors interact with the vending machine system:

\begin{itemize}
\item
\protect\phantomsection\label{bQJJXQ0cZ_VvYtiC2Ske8}
\textbf{Customer:} A primary actor who can view and select products, insert money, and collect their purchased item and any change.
\item
\protect\phantomsection\label{61hLyZTPBB6OuyGM2BERj}
\textbf{Operator:} A specialized actor with all the capabilities of a Customer, but can also perform administrative tasks like adding/removing products and collecting cash from the machine.
\end{itemize}

\subsection{Important classes and relationships}\label{1ZDLEME3nNi2PRZ7bWtNR}
\\
The design is built around several core classes that model the system's entities and their interactions. The table below summarizes these classes and their most important relationships.

\begin{center}
\begin{tabular}{|p{0.20\linewidth}|p{0.45\linewidth}|p{0.30\linewidth}|}
\hline
\textbf{Class} & \textbf{Description} & \textbf{Important Relationships} \\
\hline
VendingMachine & 
The central singleton class that manages the current state, balance, and inventory. It orchestrates the overall workflow. &
Composition with Inventory; Association with State. \\
\hline

State &
An interface defining the core user actions (insertMoney, selectProduct). &
Implemented by NoMoneyInsertedState, MoneyInsertedState, and DispenseState. \\
\hline

NoMoneyInsertedState &
A concrete state where the machine is idle and waiting for money. Only accepts money insertion. &
Implements the State interface. \\
\hline

MoneyInsertedState &
A concrete state activated after money is inserted. Handles product selection and payment validation. &
Implements the State interface. \\
\hline

DispenseState &
A concrete state responsible for dispensing the product and change, then resetting the machine. &
Implements the State interface. \\
\hline

Inventory &
Manages the collection of all racks within the machine. &
Composition with Rack. \\
\hline

Rack &
Represents a physical rack in the machine that holds a specific product and its quantity. &
Composition with Product. \\
\hline

Product &
A data class storing an item's name, price, and type. &
Associated with a Rack. \\
\hline
\end{tabular}
\end{center}


\subsection{Design highlights}\label{V7ZyAtziACDAh9e5gMRyT}
\\
This section covers the key design decisions, principles, and patterns applied to build a robust and maintainable system.

\subsubsection{Design approach}\label{CfnVVzdtw-i6xVJb2_6mw}
\\
The system was designed using a bottom-up approach. We started by identifying the simplest core entities like, \texttt{Product}, and \texttt{Rack}, then built upon them to create more complex components, like \texttt{Inventory} and the state-driven \texttt{VendingMachine} itself. This approach ensures that fundamental components are well-defined before being integrated into the larger system.

\subsubsection{Design principles}\label{o1AQTLI7nKahUF6rylcsY}
\\
The design adheres to the five SOLID principles to ensure it is robust, maintainable, and scalable.

\begin{itemize}
\item
\protect\phantomsection\label{ModuoE4y644M5n3eXRF4m}
\textbf{Single Responsibility principle (SRP):} Each class has a well-defined purpose. For example, the \texttt{Inventory} class manages racks, while the \texttt{State} classes handle state-specific behaviors.
\item
\protect\phantomsection\label{Myb_h5ClAJU55sAlgn6w5}
\textbf{Open/Closed principle (OCP):} The system is open for extension but closed for modification. Using the State pattern, we can introduce new states (e.g., a \texttt{MaintenanceState}) without altering the existing \texttt{VendingMachine} or state classes.
\item
\protect\phantomsection\label{o4yYxMdfPd9Xf82DtvX-A}
\textbf{Liskov Substitution principle (LSP):} All concrete state classes (\texttt{NoMoneyInsertedState}, \texttt{MoneyInsertedState}, \texttt{DispenseState}) are substitutable for the \texttt{State} interface. The \texttt{VendingMachine} interacts with them through the common interface, ensuring consistent behavior.
\item
\protect\phantomsection\label{XZAU3OeV6CpKQ8VWYKKZk}
\textbf{Interface Segregation principle (ISP):} The \texttt{State} interface is lean and specific to user actions. Administrative functions like \texttt{addRack} are kept within the \texttt{VendingMachine} class, so state classes are not forced to implement methods they don't need.
\item
\protect\phantomsection\label{fMdBrcJ-4BHFBhPkkl6Qn}
\textbf{Dependency Inversion principle (DIP):} The \texttt{VendingMachine} class depends on the \texttt{State} interface (an abstraction) rather than its concrete implementations. This decouples the main class from the details of state-specific logic.
\end{itemize}

\subsubsection{Design patterns}\label{GlqRsR5Gb4uRmWqxe71rO}
\\
The design leverages established design patterns to solve common problems effectively.

\begin{itemize}
\item
\protect\phantomsection\label{rdYrMl-173UglOP7ZxA4e}
\textbf{State pattern:} This pattern is central to the design. It encapsulates the vending machine's varying behavior into separate state objects (\texttt{NoMoneyInsertedState}, \texttt{MoneyInsertedState}, \texttt{DispenseState}), allowing it to alter its behavior when its internal state changes. This avoids complex conditional logic within the main \texttt{VendingMachine} class.
\end{itemize}

\subsection{Object interactions}\label{1w9GeQp10VXbtONmzs_bo}
\\
This section describes how key objects collaborate to fulfill user requests, focusing on the flow of method calls between them.

\begin{itemize}
\item
\protect\phantomsection\label{ZpE5R3XJRPYoebXpipJDd}
\textbf{Money insertion:} When a \texttt{Customer} inserts money, the \texttt{VendingMachine} delegates the call to its current state object. If the state is \texttt{NoMoneyInsertedState}, it accepts the money, updates the machine's balance, and transitions the machine's state to \texttt{MoneyInsertedState}.
\item
\protect\phantomsection\label{rHRauEx0XJOgUJ69mavVr}
\textbf{Product selection:} A \texttt{Customer} selects a product by its rack number. The \texttt{VendingMachine} passes this request to the current \texttt{MoneyInsertedState} object. The state object then collaborates with the \texttt{Inventory} to check the product's availability and price.
\item
\protect\phantomsection\label{26KgRF8xEvG6bXzD3Gaj1}
\textbf{Payment validation and dispensing:} The \texttt{MoneyInsertedState} object compares the machine\textquotesingle s current balance to the product's price.

\begin{itemize}
\item
If payment is sufficient, it instructs the \texttt{VendingMachine} to dispense the product and any required change. The machine's state then transitions to \texttt{DispenseState}.
\item
If payment is insufficient, it instructs the \texttt{VendingMachine} to refund the money and transitions the state back to \texttt{NoMoneyInsertedState}.
\end{itemize}
\item
\protect\phantomsection\label{JUfoJqmOBgt861CLVj-KS}
\textbf{Dispensing and reset:} Once in \texttt{DispenseState}, the machine dispenses the item and change. After the process, the state automatically returns to \texttt{NoMoneyInsertedState} to prepare for the next transaction.
\end{itemize}

\subsection{System workflow}\label{qvlWcAX1I_sE1WprJZ_9K}
\\
This section outlines the primary end-to-end scenario of a customer purchasing a product, as visualized in the activity diagram.

\begin{itemize}
\item
\protect\phantomsection\label{dUU1udyYwOxFXCGCK61au}
\textbf{Initial state:} The process begins with the machine idle, waiting for a user to insert money.
\item
\protect\phantomsection\label{U8q5ahRT_oIMqS187EAlk}
\textbf{Money insertion:} A customer inserts cash into the machine.
\item
\protect\phantomsection\label{xxas9Q4WtmZSOh37uTZVi}
\textbf{Product selection:} The customer selects a desired product. The system then checks if the selected product is available in the inventory. If not, it informs the customer and prompts them to make another selection.
\item
\protect\phantomsection\label{JQ6UF8eAW7t5ukCkN01nb}
\textbf{Payment validation:} If the product is available, the system validates the payment.

\begin{itemize}
\item
If the inserted amount exceeds the product's price, the machine dispenses the money back to the customer, and the flow returns to the money insertion step.
\item
If the amount is exactly the price, the system dispenses the product.
\item
If the amount exceeds the price, the system first dispenses the correct change and then dispenses the product.
\end{itemize}
\item
\protect\phantomsection\label{ZznUgd7VFw_I83urjEvzj}
\textbf{Final state:} The customer receives their product (and any change), concluding the transaction. The machine then resets for the next customer.
\end{itemize}

% End of section content